\documentclass{article}

\usepackage[T1]{fontenc}
\usepackage{textcomp}
\usepackage{microtype}
\usepackage{graphicx}
\usepackage{subfigure}
\usepackage{booktabs}

\usepackage{hyperref}

\newcommand{\theHalgorithm}{\arabic{algorithm}}

\usepackage[accepted]{icml2026}

\icmltitlerunning{A Compact and Fast Model for Vietnamese Diacritic Restoration}

\begin{document}

\twocolumn[
\icmltitle{A Compact, Fast Model for Vietnamese Diacritic Restoration}

{\centering\textbf{Project Category:} Natural Language\par}

\begin{icmlauthorlist}
\icmlauthor{\begin{tabular}{@{}c@{}}\textbf{Pham Tuan Hung}\\[1pt]\small 24hung.pt@vinuni.edu.vn\end{tabular}}{vinuni}
\icmlauthor{\begin{tabular}{@{}c@{}}\textbf{Do Quang Thai An}\\[1pt]\small 24an.dqt@vinuni.edu.vn\end{tabular}}{vinuni}
\icmlauthor{\begin{tabular}{@{}c@{}}\textbf{Nguyen Ha Bao Vy}\\[1pt]\small 24vy.nhb@vinuni.edu.vn\end{tabular}}{vinuni}
\end{icmlauthorlist}

\icmlaffiliation{vinuni}{VinUniversity, Hanoi, Vietnam}

\icmlcorrespondingauthor{Pham Tuan Hung}{24hung.pt@vinuni.edu.vn}

\icmlkeywords{Vietnamese diacritic restoration, natural language processing, efficiency}

\vskip 0.1in
]

\printAffiliationsAndNotice{}

\begin{abstract}
We study Vietnamese diacritic restoration: recovering the tone marks and
diacritics of text written without them. Existing systems are accurate but large and
slow, and restoration runs on every sentence. We take the
opposite end of the trade-off: a small CPU model where a factorized character-level
network scores only the spellings attested for each form, so an invalid syllable
cannot be produced. We evaluate on multi-register Vietnamese text from public
sources, reporting token and sentence accuracy, parameter count, and CPU latency.
\end{abstract}

\section{Introduction}
\label{sec:intro}

\textbf{Motivation.}
Much Vietnamese text is written without diacritics: chat messages, comments,
search queries, speech output, and digitised legacy data. Restoring them is hard
because one undiacritised string maps to many words; \textit{ban} may be
\textit{b\d{a}n} (friend), \textit{b\`an} (table), or \textit{b\'an} (to sell). Correct resolution needs lexical knowledge and sentence context. Since
restoration runs on every sentence, often on a phone or CPU server, the goal is
not accuracy alone but accuracy at a small fixed cost. This is an
\emph{application project}: we build a model that is small and fast without giving
up accuracy.

\textbf{Limitations of existing approaches.}
Character-by-character models reach accuracy in the high nineties
\citep{do2013mt,nga2019deep}, but they can emit sequences that are not valid
Vietnamese syllables, and beam search adds latency for a step run on every
sentence \citep{dang2020tdp}. Current alternatives are large
sequence-to-sequence or language models, whose per-sentence cost is orders of
magnitude above a small local network. Existing work reports accuracy but rarely
reports it \emph{with} parameter count and CPU latency on the same data.

\textbf{Method.}
We apply supervised classification over a restricted label space. A
character-level neural encoder turns each undiacritised syllable into a vector,
and two classification heads sit on top: one over the six tone classes, the other
over vowel modifications. This factorises one large decision into two small
classifications. Unlike earlier systems that predict characters freely, both heads
score only the spellings attested for the input form, so the label space is
exactly the set of real Vietnamese syllables. Training minimises cross-entropy
with AdamW by backpropagation. Since decoding is a single argmax, no beam search
is required, and the network stays small enough for CPU.

\textbf{Data.}
Diacritised Vietnamese is abundant, and training pairs are built without human
labelling by stripping diacritics. We draw on Vietnamese Wikipedia, open news corpora
(BKAINewsCorpus, vietnamnet-news) and further open sources in other registers,
since Wikipedia under-represents conversational and literary text. Sources are
split and deduplicated at article level before sampling, so no sentence crosses
the train/validation boundary.

\textbf{Intended experiments.}
We report token and exact-sentence accuracy per register with parameter count
and CPU latency per sentence, so accuracy is read against cost. Errors are
stratified by candidate-set size to separate ambiguous from unambiguous syllables,
and we compare against a frequency baseline picking the commonest spelling of each
form, showing what is learned beyond vocabulary.

\bibliography{references}
\bibliographystyle{icml2025}

\end{document}
